\SourceLeaf{219}%
\noindent nach der Annahme gilt, dass aber der Satz

\[ c\;[\text{1. Figur}] = c\;[\text{2. Figur}] \]

\noindent dann bewiesen werden müsse?

Da die Winkel beliebig sind, scheint mir dieser Beweis nicht notwendig.

Dies sind die Grundzüge des Beweises, und ich erwarte Ihre Entscheidung. Zur Rechtfertigung meines Beweises füge ich nur hinzu, dass durch den zweiten Schritt zwar das Dreieck $ABC$ verschwindet, aber \textls{nicht} die Winkel des Dreiecks. Wie die Linien auch liegen, es gilt stets

\[ J\hat{B}H = \beta, \qquad G\hat{C}F = \gamma, \qquad D\hat{A}E = \alpha \]

\noindent im endlichen ebenso wie im verschwindenden Dreieck. Somit ist die Summe

\[ J\hat{A}H + G\hat{A}F + D\hat{A}E \]

\noindent immer gleich der Summe der Winkel eines geradlinigen Dreiecks.

Soll man also den Satz für ein beliebiges Dreieck (mit den Winkeln $A$, $B$, $C$) beweisen, so zieht man die Linien $DG$, $EH$, sodass

\[ \alpha = A, \]

\noindent und nimmt ferner den Winkel $J\hat{A}H = B$, $G\hat{A}F = C$.

Ist dann $JAF$ keine gerade, sondern eine gebrochene Linie $JAF'$, so ist der Winkel $c$ dadurch zwar um $dc$ kleiner, der Winkel $b$ aber um ebenso viel größer geworden. Somit bleibt ihre Summe unverändert, oder wir haben, was für die Stringenz des Beweises erforderlich ist:

\[ b + c\;[\text{Fig. 1}] = b + c\;[\text{Fig. 2}].\} \]

\bigskip
\hrule
\bigskip

\Needspace{7\baselineskip}
\begin{center}
\textsc{Gauss} an \textsc{Schumacher}. Göttingen, 17. Mai 1831.
\end{center}

\mbox{.\ .\ .\ .\ .\ .\ .} Bei dem, was Sie über die Parallellinien schreiben, haben Sie bei genauer Betrachtung in Ihren Syllogismen einen Zwischensatz gebraucht, ohne ihn ausdrücklich auszusprechen. Er müsste so lauten:
